% answers_short.tex —— 2012 年·全国大纲卷·文科 · 简短答案速查
% 由 tools/make_short_answers.py 自 answers.tex 的 \gksolhead/\gkanswer 提取，
% 并与 2012/2012_数学真题_全国大纲文_答案要点.md 逐题比对后生成，勿手改（改源 MD 后重跑）。

\begin{gknotice}
\textbf{使用说明}：本节只给最终答案便于快速核对，完整推导见解析版。答案为编写者独立推导，\textbf{非官方标准答案}。
\end{gknotice}

\gkpart{选择题}
\begin{center}\begin{tabular}{ccccccccccccc}
\toprule
题号 & 1 & 2 & 3 & 4 & 5 & 6 & 7 & 8 & 9 & 10 & 11 & 12 \\
\midrule
答案 & B & A & C & A & C & B & C & D & D & C & D & B \\
\bottomrule
\end{tabular}\end{center}

\gkpart{填空题}
\begin{itemize}[leftmargin=2.6em]
  \item[\textbf{13.}] $7$
  \item[\textbf{14.}] $-1$
  \item[\textbf{15.}] $\dfrac{5\pi}{6}$
  \item[\textbf{16.}] $\dfrac35$
\end{itemize}

\gkpart{解答题}
\begin{itemize}[leftmargin=2.6em]
  \item[\textbf{17.}] $A=30^\circ$ 或 $A=90^\circ$。
  \item[\textbf{18.}] (1) $a_2=3$，$a_3=6$；(2) $a_n=\dfrac{n(n+1)}2$（$n\in\mathbb N^+$）
  \item[\textbf{19.}] (1) 见解析版；(2) $PD$ 与平面 $PBC$ 所成角为 $30^\circ$
  \item[\textbf{20.}] (1) $\dfrac{36}{125}$；(2) $\dfrac{192}{625}$
  \item[\textbf{21.}] (1) $a\geq0$ 时 $f$ 在 $\mathbb R$ 上单调递增；$a<0$ 时，令 $s=\sqrt{-a}$，则 $f$ 在 $(-\infty,-1-s)$、$(-1+s,+\infty)$ 上递增，在 $(-1-s,-1+s)$ 上递减；(2) $a=-1$
  \item[\textbf{22.}] (1) $r=\dfrac{\sqrt5}{2}$；(2) $\dfrac{6\sqrt5}{5}$
\end{itemize}
